\documentclass[a4 paper,fontset=windows]{ctexart}
\usepackage{fontspec}
\usepackage{color}
\definecolor{dkgreen}{rgb}{0,0.5,0}
\definecolor{dkred}{rgb}{0.5,0,0}
\usepackage{graphicx}
\setsansfont{Consolas}
\setmonofont{Consolas}
\usepackage{geometry}
\geometry{top=2.5cm,bottom=2.5cm,left=2.5cm,right=2.5cm}

\usepackage{listings}
\lstset{
	language=Python,
	basicstyle=\ttfamily\small,
	aboveskip={1.0\baselineskip},
	belowskip={1.0\baselineskip},
	columns=fixed,
	extendedchars=true,
	breaklines=true,
	tabsize=4,
	prebreak=\raisebox{0ex}[0ex][0ex]{\ensuremath{\hookleftarrow}},
	frame=lines,
	showtabs=false,
	showspaces=false,
	showstringspaces=false,
	keywordstyle=\color[rgb]{0,0.5,0},
	commentstyle=\color[rgb]{0.5,0,0},
	stringstyle=\color[rgb]{1,0,0},
	numbers=left,
	numberstyle=\small\ttfamily ,
	stepnumber=1,
	numbersep=12pt,
	captionpos=t,
	escapeinside={\%*}{*)}
	}

\begin{document}
	
	\title{\textbf{数据结构与算法K05\ \ \ \ 课堂练习}}
	\author{张植竣\ 1900011604\ 北京大学物理学院天文学系}
	\date{2020年3月24日}
	\maketitle
	
	\section*{绘制黄金螺线}
	
	代码如下：

	\begin{lstlisting}[caption = Golden Curve]
	import turtle as t

	def f(x):  #斐波那契数列的第x项
		if x==1 or x==2:
			return 1
		else:
			return f(x-1)+f(x-2)

	def rot(x):  #n阶卷
		if x >= 1:
			t.circle(4 * f(x), 90)
			rot(x-1)
		else:
			t.hideturtle()
			t.done()
	
	rot(10)
	\end{lstlisting}
	
	输出图片如下：
	
	\begin{figure}[ht]
		\centering
		\includegraphics[width=0.32\linewidth]{1900011604_K05_picture.pdf}
	\end{figure}

\end{document}